\section{The Real and Complex Number Systems}
\subsection{Introduction}
\begin{example}
It is well known that there is no rational $p$ such that $p^2=2$. 

Define $A \coloneqq \left\{p \in \Q^+ \mid p^2<2\right\}$ and $B \coloneqq \left\{p \in \Q^+ \mid p^2>2\right\}$. 
Then $A$ contains no largest number and $B$ contains no smallest.
\end{example}
\begin{proof}
It suffices to show that
\[
\forall p \in A, \, \exists q \in A \, s.t. \, p<q
\]
and 
\[
\forall p \in B, \, \exists q \in B \, s.t. \, q<p \mathrlap{.}
\]

For every $p \in \Q^+$, take 
\begin{equation}\label{eq:q}
    q = p - \frac{p^2-2}{p+2} = \frac{2p+2}{p+2}.
\end{equation}
Then
\begin{equation}\label{eq:q^2-2}
q^2-2 = \frac{2(p^2-2)}{(p+2)^2}.
\end{equation}


Hence,
\begin{align*}
p \in A &\implies p^2-2<0 \\
&\implies p<q \text{ and } q^2<2 \tag*{$\because \eqref{eq:q}, \eqref{eq:q^2-2}$}\\
&\implies q \in A
\shortintertext{and}
p \in B &\implies p^2-2>0  \\
&\implies 0<q<p \text{ and } q^2>2  \tag*{$\because \eqref{eq:q}, \eqref{eq:q^2-2}$}\\
&\implies q \in B. 
\end{align*}


We have shown the desired result.
\end{proof}

\subsection{Ordered Sets}
\begin{definition}
Let $S$ be a set. An \emph{order} on $S$ is a relation, denoted by $<$, with the following two properties:
\begin{enumerate}
\item[(i)] If $x \in S$ and $y \in S$ then one and only one of the statements 
\[
x<y, \qquad x=y, \qquad y<x
\]
is true.
\item[(ii)] If $x,y,z \in S$, if $x<y$ and $y<z$, then $x<z$.
\end{enumerate}
\end{definition}
\begin{definition}
An \emph{ordered set} is a set $S$ in which an order is defined.
\end{definition}
\begin{definition}
Suppose $S$ is an ordered set, and $E \subset S$. If there exists a $\beta \in S$ such that $x \leq \beta$ for every $x \in E$, we say that $E$ is \emph{bounded above}, and call $\beta$ an \emph{upper bound} of $E$.

Lower bounds are defined in the same way (with $\geq$ in place of $\leq$).
\end{definition}
\begin{definition}
Suppose $S$ is an ordered set, $E \subset S$, and $E$ is bounded above. Suppose there exists an $\alpha \in S$ with the following properties: 
\begin{enumerate}
    \item[(i)] $\alpha$ is an upper bound of $E$.
    \item[(ii)] If $\gamma < \alpha$ then $\gamma$ is not an upper bound of $E$.
\end{enumerate}

Then $\alpha$ is called the \emph{least upper bound of $E$} [that there is at most one such $\alpha$ is clear from (ii)] or the \emph{supremum of $E$}, and we write \[\alpha=\sup E.\]

The \emph{greatest lower bound}, or \emph{infimum}, of a set $E$ which is bounded below is defined in the same manner: The statement \[\alpha=\inf E\] means that $\alpha$ is a lower bound of $E$ and that no $\beta$ with $\beta > \alpha$ is a lower bound of $E$.
\end{definition}
\begin{definition}
An ordered set $S$ is said to have the \emph{least-upper-bound property} if the following is true:

If $E \subset S$, $E$ is not empty, and $E$ is bounded above, then $\sup E$ exists in $S$.
\end{definition}
\begin{remark}
    The ordered set $\Q$ does not have the least-upper-bound property.
\end{remark}

The following theorem shows that every ordered set with the least-upper-bound property also has the greatest-lower-bound property.
\begin{theorem}
Suppose $S$ is an ordered set with the least-upper-bound property, $B \subset S$, $B$ is not empty, and $B$ is bounded below. Let $L$ be the set of all lower bounds of $B$. Then \[\alpha = \sup L\] exists in $S$, and $\alpha = \inf B$.

In particular, $\inf B$ exists in $S$.
\end{theorem}
\begin{proof}
Since $B$ is bounded below, $L$ is nonempty. 
Let $x \in B$. Then $y \leq x$ for every $y \in L$. 
Thus, every $x \in B$ is an upper bound of $L$. 
Since $L$ is a nonempty subset of $S$ that is bounded above, by hypothesis, \[\alpha = \sup L\] exists in $S$.

If $\gamma < \alpha$ then $\gamma$ is not an upper bound of $L$, hence $\gamma \notin B$. It follows that $\alpha \leq x$ for every $x \in B$. Thus, $\alpha \in L$.

If $\alpha < \beta$ then $\beta \notin L$, since $\alpha$ is an upper bound of $L$.

We conclude that $\alpha$ is a lower bound of $B$ (since $\alpha \in L$) and if $\alpha < \beta$ then $\beta$ is not a lower bound of $B$ (since $\beta \notin L$), which implies $\alpha = \inf B$.
\end{proof}
\begin{remark}
Fix $x \in B$. Then $y \leq x$ for all $y \in L$. Thus, $\sup L \leq x$. Since the choice of $x \in B$ was arbitrary, we have 
\[
\sup L = \alpha \leq x \text{ for every } x \in B,
\]
which implies $\alpha$ is a lower bound of $B$. Hence, $\alpha \in L$.
\end{remark}
\subsection{Fields}
\begin{definition}
A \emph{field} is a set $F$ with two operations, called \emph{addition} and \emph{multiplication}, which satisfy the following so-called "field axioms" (A), (M), and (D):
\begin{enumerate}
\item[(A)] Axioms for addition
\begin{itemize}
    \item[(A1)] If $x \in F$ and $y \in F$, then their sum $x+y$ is in $F$.
    \item[(A2)] Addition is commutative: $x+y=y+x$ for all $x,y \in F$.
    \item[(A3)] Addition is associative: $(x+y)+z=x+(y+z)$ for all $x,y,z \in F$.
    \item[(A4)] $F$ contains an element $0$ such that $0+x=x$ for every $x \in F$.
    \item[(A5)] To every $x \in F$ corresponds an element $-x \in F$ such that \[x+(-x)=0.\]
\end{itemize}
\item[(M)] Axioms for multiplication
\begin{itemize}
    \item[(M1)] If $x \in F$ and $y \in F$, then their product $xy$ is in $F$.
    \item[(M2)] Multiplication is commutative: $xy=yx$ for all $x,y \in F$.
    \item[(M3)] Multiplication is associative: $(xy)z=x(yz)$ for all $x,y,z \in F$.
    \item[(M4)] $F$ contains an element $1 \neq 0$ such that $1x=x$ for every $x \in F$.
    \item[(M5)] If $x \in F$ and $x \neq 0$ then there exists an element $1/x \in F$ such that \[x \cdot (1/x) = 1.\]
\end{itemize}
\item[(D)] The distributive law
\[
x(y+z)=xy+xz
\]
holds for all $x,y,z \in F$.
\end{enumerate}
\end{definition}
\begin{definition}
An \emph{ordered field} is a \emph{field} $F$ which is also an \emph{ordered set}, such that
\begin{itemize}
    \item[(i)] $x+y<x+z$ if $x,y,z \in F$ and $y<z$.
    \item[(ii)] $xy>0$ if $x \in F$, $y \in F$, $x>0$, and $y>0$.
\end{itemize}
\end{definition}

\subsection{The Real Field}
\begin{theorem}
There exists an ordered field $\R$ which has the least-upper-bound property.

Moreover, $\R$ contains $\Q$ as a subfield.
\end{theorem}

\begin{theorem}
    \item[(a)] If $x \in \R$, $y \in \R$, and $x>0$, then there is a positive integer $n$ such that \[nx>y.\]
    \item[(b)] If $x \in \R$, $y \in \R$, and $x<y$, then there exists a $p \in \Q$ such that $x<p<y$.
\end{theorem}
\begin{proof}
\item[(a)] 
Let $x,y \in \R$ and $x>0$. 
Suppose to the contrary that $nx \leq y$ for each $n \in \N$ and define \[A \coloneqq \left\{nx \mid n \in \N \right\}.\] 
Then $A$ is nonempty and bounded above by $y$. 
Thus, the supremum of $A$ exists in $\R$. Put $\alpha = \sup A$.

Since $x>0$, we have $\alpha - x < \alpha$, hence, $\alpha-x$ is not an upper bound of $A$. Therefore, $\alpha-x<mx$ for some $m \in \N$. 
Then $\alpha < (m+1)x \in A$, which is a contradiction, since $\alpha$ is an upper bound of $A$.
\item[(b)]
Since Rudin's proof of part (b) is already explicit, we omit it here.
\end{proof}
\begin{theorem}\label{thm:nth_square_root}
For every real $x>0$ and every integer $n>0$ there is one and only one positive real $y$ such that $y^n=x$.
\end{theorem}
This number $y$ is written $\sqrt[n]{x}$ or $x^{1/n}$. 
\begin{proof}
We omit the proof here.
\end{proof}
\begin{corollary}\label{thm:associative_nth_sqrt}
If $a$ and $b$ are positive and real numbers and $n$ is a positive integer, then 
\[
(ab)^{\frac{1}{n}}=a^{\frac{1}{n}}b^{\frac{1}{n}}.
\]
\end{corollary}
\begin{proof}
We omit the proof here.
\end{proof}
\begin{lemma}\label{greatest integer}
\item[(a)] If $E \subset \Z$ has the infimum, then $\inf E \in E$.
\item[(b)] If $x \in \R$ then there exists an $n \in \Z$ such that \[n \leq x < n+1.\]
\end{lemma}
\begin{proof}
\item[(a)] Let $\alpha = \inf E$.
Suppose to the contrary that $\alpha \notin E$. 
Since $\alpha+1$ is not a lower bound of $E$, choose an $x_0 \in E$ such that 
\begin{equation}\label{ineq1}
    \alpha<x_0<\alpha+1.
\end{equation}
Since $x_0$ is not a lower bound of $E$, choose an $x_1 \in E$ again such that
\begin{equation}\label{ineq2}
    \alpha<x_1<x_0.
\end{equation}

Subtracting $x_0$ from \eqref{ineq2}, we obtain $\alpha-x_0<x_1-x_0<0$. 
Since $-1<\alpha-x_0$ from \eqref{ineq1}, it follows that $-1<x_1-x_0<0$.
This contradicts the fact that $x_1-x_0 \in \Z$.
We conclude that $\alpha \in E$.
\item[(b)] Let $x \in \R$ and define $E \coloneqq \left\{n \in \Z \mid n>x \right\}$.
By the Archimedean Property, $E$ is nonempty. 
Since $E$ is bounded below by $x$, the infimum of $E$ exists in $\R$. 
Since $E \subset \Z$, $\inf E$ is the least element of $E$ by (a).
Thus, we have \[\inf E -1 \leq x < \inf E.\]
Putting $n= \inf E -1 \in \Z$, we obtain the desired inequality.
\end{proof}
\begin{remark}[Decimals]
Let $x>0$ be real. 
Let $n_0$ be the largest integer such that $n_0 \leq x$, by lemma \ref{greatest integer}.
Having chosen $n_0, n_1, \ldots, n_{k-1}$, let $n_k$ be the largest integer such that 
\[
n_0 + \frac{n_1}{10} + \cdots + \frac{n_k}{10^k} \leq x.
\]
Let $E$ be the set of these numbers 
\begin{equation} \label{dec1}
n_0 + \frac{n_1}{10} + \cdots + \frac{n_k}{10^k} \qquad (k=0, 1, 2, \ldots).
\end{equation}
Notice that $x$ is an upper bound of $E$ and if $k \in \N$ then
\[
n_0 + \frac{n_1}{10} + \cdots + \frac{n_k}{10^k} \leq x < n_0 + \frac{n_1}{10} + \cdots + \frac{n_k+1}{10^k}.
\]
Hence, we have
\begin{equation} \label{ineq:error_term}
x - n_0 - \frac{n_1}{10} - \cdots - \frac{n_k}{10^k} < \frac{1}{10^k} \text{ for } k  \in \N.
\end{equation}

Suppose $y<x$. Then by the Archimedean Property, there exists a positive integer $k$ such that $1/10^k<x-y$.
From \eqref{ineq:error_term}, we obtain 
\[
x - n_0 - \frac{n_1}{10} - \cdots - \frac{n_k}{10^k}<x-y,
\]
hence 
\[
y<n_0 + \frac{n_1}{10} + \cdots + \frac{n_k}{10^k}.
\]
Therefore, $x=\sup E$. The decimal expansion of $x$ is 
\begin{equation} \label{dec2}
n_0 \cdot n_1 n_2 n_3 \cdots .
\end{equation}

Conversely, for any infinite decimal \eqref{dec2}, the set $E$ of numbers \eqref{dec1} is bounded above, and \eqref{dec2} is the decimal expansion of $\sup E$.
\end{remark}
\subsection{Exercises}
\begin{exercise}
Let $A$ be a nonempty set of real numbers which is bounded below.
Let $-A$ be the set of all numbers $-x$, where $x \in A$.
Prove that \[\inf A = -\sup (-A).\]
\end{exercise}
\begin{proof}
Since $A$ is bounded below, the infimum of $A$ exists in $\R$. 
If $a \in A$ then $\inf A \leq a$. Hence,
\begin{equation}\label{ineq:ex_1_a}
-a \leq -\inf A \text{ for all } a \in A,
\end{equation}
and it follows that $-A$ is bounded above. Thus, $\sup (-A)$ exists in $\R$.
Since $- \inf A$ is an upper bound of $-A$ from \eqref{ineq:ex_1_a}, we have $\sup (-A) \leq -\inf A$, hence 
\begin{equation}\label{ineq:ex_1_b}
\inf A \leq -\sup(-A).
\end{equation}
Since $-a \leq \sup(-A)$ for every $a \in A$, we obtain $-\sup (-A) \leq a$ for all $a \in A$.
Therefore, $-\sup(-A)$ is a lower bound of $A$ and 
\begin{equation}\label{ineq:ex_1_c}
-\sup(-A) \leq  \inf A.
\end{equation}
Combining \eqref{ineq:ex_1_b} and \eqref{ineq:ex_1_c}, we conclude that $\inf A = -\sup (-A)$.
\end{proof}
\begin{exercise}
Fix $b>1$.
\item[(a)] If $m$, $n$, $p$, $q$ are integers, $n>0$, $q>0$, and $r=m/n=p/q$, prove that 
 \[
 (b^m)^{1/n}=(b^p)^{1/q}.
 \]
 Hence it makes sense to define $b^r = (b^m)^{1/n}$.
 \item[(b)] Prove that $b^{r+s} = b^r b^s$ if $r$ and $s$ are rational.
 \item[(c)] If $x$ is real, define $B(x)$ to be the set of all numbers $b^t$, where $t$ is rational and $t \leq x$. 
Prove that \[b^r = \sup B(r)\] when $r$ is rational.
Hence it makes sense to define \[b^x = \sup B(x)\] for every real $x$.
\item[(d)] Prove that $b^{x+y}=b^xb^y$ for all real $x$ and $y$.

\end{exercise}
\begin{proof}
\item[(a)]
Let $b^{mq}=b^{np}=x$.
Since $n$ and $q$ are positive integers and $b>0$, by Theorem \ref{thm:nth_square_root}, we have 
\[
b^m=x^{1/q} \text{ and } b^p=x^{1/n}.
\]
It follows that
\[
(b^m)^{1/n}=(x^{1/q})^{1/n}=(x^{1/n})^{1/q}=(b^p)^{1/q},
\]
since $[(x^{1/q})^{1/n}]^{nq}=[(x^{1/n})^{1/q}]^{nq}=x$.
\item[(b)]Let $m$, $n$, $p$, $q$ be integers and $n>0$, $q>0$ and let $r=m/n$ and $s=p/q$. 
Then by Corollary \ref{thm:associative_nth_sqrt}, we obtain
\begin{align*}
b^{r+s}
&=b^{\frac{m}{n}+\frac{p}{q}}
=b^{\frac{mq+np}{nq}}\\
&=\left(b^{mq+np}\right)^{\frac{1}{nq}}\\
&=\left(b^{mq}b^{np}\right)^{\frac{1}{nq}}\\
&=\left(b^{mq}\right)^{\frac{1}{nq}} \left(b^{np}\right)^{\frac{1}{nq}}\\
&=b^{\frac{mq}{nq}} b^{\frac{np}{nq}}
=b^{\frac{m}{n}} b^{\frac{p}{q}}
=b^rb^s.
\end{align*}
\item[(c)]
Note that by Theorem \ref{thm:nth_square_root}, $b^r>0$ if $b>1$ and $r$ is rational.

Let $n$ be a positive integer and let $0<a<1$ be real. We first show that $0<a^{1/n}<1$.

Suppose to the contrary that $a^{1/n} \geq 1$. Then we obtain 
\[
1 \leq \left(a^{1/n}\right)^n=a<1,
\]
a contradiction. Therefore, if $0<a<1$ then $0<a^{1/n}<1$, where $n$ is a positive integer.

Suppose $r$ and $s$ are rational. 
Let $m$, $n$, $p$, and $q$ be integers and let $n$ and $q$ be positive.
Then we write $r=m/n$, $s=p/q$.

Suppose $r<s$. Then
\[
r-s =\frac{m}{n}-\frac{p}{q} =\frac{mq-np}{nq} <0,
\]
hence
\[
0 < b^{mq-np} = \frac{1}{b^{np-mq}} < 1
\]
and 
\[
0 < b^{r-s} = b^{\frac{mq-np}{nq}} = \left(b^{mq-np}\right)^{1/nq} < 1.
\]
Thus, if $r<s$ then
\[
b^r = b^{(r-s)+s} = b^{r-s}b^s < b^s.
\]

If $t$ is rational and $t \leq r$ then $b^t \leq b^r$. 
It follows that $B(r)$ is bounded above by $b^r$ and $\sup B(r)$ exists in $\R$.
Since $b^r$ is an upper bound of $B(r)$, 
\[
\sup B(r) \leq b^r.
\]
Since $b^r \in B(r)$,
\[
b^r \leq \sup B(r).
\]

We conclude that $b^r = \sup B(r)$.

\item[(d)]
Let $t_1$ and $t_2$ be rational. 
If $t_1 \leq x$ and $t_2 \leq y$ then $t_1+t_2 \leq x+y$.
Thus, we have
\[
b^{t_1}b^{t_2} = b^{t_1+t_2} \leq \sup B(x+y) =b^{x+y}.
\]

Fix $t_1\leq x$. Then 
\[
b^{t_2} \leq b^{x+y}/b^{t_1} \text{ for all } t_2 \leq y.
\]
Hence, 
\[
b^y = \sup B(y) \leq b^{x+y}/b^{t_1}.
\]
Since the choice of $t_1$ was arbitrary,
\[
b^{t_1} \leq b^{x+y}/b^y \text{ for all } t_1 \leq x,
\]
and it follows that
\[
b^x = \sup B(x) \leq b^{x+y}/b^y.
\]
Therefore, we obtain
\begin{equation} \label{ineq:Ch.1_ex2_(d)_1}
    b^xb^y \leq b^{x+y}.
\end{equation}

Notice that 
\[
0 < b^{t_1} \leq b^x \text{ for every } t_1 \leq x
\]
and 
\[
0 < b^{t_2} \leq b^y \text{ for every } t_2 \leq y \mathrlap{.}
\]
Hence, for every $t_1 \leq x$ and $t_2 \leq y$, 
\[
b^{t_1+t_2} = b^{t_1}b^{t_2} \leq b^xb^y.
\]

Now, to obtain
\begin{equation}\label{ineq:Ch.1_ex2_(d)_2}
b^{x+y} = \sup B(x+y) \leq b^xb^y
\end{equation}
from above inequality, we need the following statement:
\[
b^x=\sup \left\{b^t \mid t \in \Q, t<x\right\}.
\]

Let us prove it. Notice that
\[
b^n-1=(b-1)(b^{n-1}+ \cdots + 1) \geq n(b-1)
\]
for any positive integer $n$. Put $c=b^{1/n}>1$. Then we have
\[
b-1 \geq n(b^{1/n}-1).
\]
Hence, if $t>1$ and $n>(b-1)/(t-1)$ then
\[
\frac{n(b^{1/n}-1)}{t-1} \leq \frac{b-1}{t-1} < n,
\]
and it follows that
\[
b^{1/n}<t.
\]

Note that if $t>1$ then $b^{1/n}<t$ for sufficiently large $n$.

Let $B_0(x)=\left\{b^t \mid t \in \Q, t<x\right\}$.
We need to show that $b^x$ is the supremum of $B_0(x)$. 
It is clear that $b^x$ is an upper bound of $B_0(x)$.
Thus, it remains to prove that if $y<b^x$ then $y$ is not an upper bound.

Suppose $0<y<b^x$. Then $b^x/y >1$ and it follows that
\[
b^{1/n} <b^x/y
\]
for sufficiently large $n$. 
Hence, we have
\[
y<\frac{b^x}{b^{1/n}}<b^x.
\]
By the density of rationals, choose a rational $t$ such that $x-(1/n)<t<x$. 
Since $x<s<t+(1/n)$ for some rational $s$, 
\[
b^x \leq b^s <b^{t+(1/n)} = b^tb^{1/n}
\]
and
\[
y<\frac{b^x}{b^{1/n}}<b^t.
\]
Thus, $y$ is not an upper bound of $B_0(x)$ and therefore $b^x = \sup B_0(x)$.

Now, let $t$ be a rational and let $t<x+y$.
Pick $t_2 \in \Q$ such that $t-x<t_2<y$ and set $t_1=t-t_2$.
Then for every $t \in \Q$ with $t<x+y$, there are rational $t_1<x$ and $t_2<y$ such that $t=t_1+t_2<x+y$, hence $b^t=b^{t_1+t_2}=b^{t_1}b^{t_2} \leq b^xb^y$.
Thus, we obtain \eqref{ineq:Ch.1_ex2_(d)_2}.

Combining \eqref{ineq:Ch.1_ex2_(d)_1} and \eqref{ineq:Ch.1_ex2_(d)_2}, we conclude that
\[
b^{x+y}=b^xb^y.
\]
\end{proof}
\begin{exercise}
Prove that no order can be defined in the complex field that turns it into an ordered field.
\end{exercise}
\begin{proof}
Suppose to the contrary that there exists an order $<$ on $\C$ which turns $\C$ into an ordered field.
Let $\mathbf{0}$ and $\mathbf{1}$ denote the additive and multiplicative identities of the complex field, respectively. 
Then $\mathbf{0}=(0,0)$ and $\mathbf{1}=(1,0)$. 

Note that $i = (0,1) \neq \mathbf{0}$.
If $x \neq \mathbf{0}$ then $x^2 > \mathbf{0}$. 
It follows that
\begin{align*}
    i^2=(0,1)(0,1)=(-1,0)>\mathbf{0}.
\end{align*}
Since $(-1,0)+(1,0)=\mathbf{0}$, $(-1,0)$ is the additive inverse of $(1,0)$. 
Since $\mathbf{1}>\mathbf{0}$ and if $x>\mathbf{0}$ then $-x<\mathbf{0}$ in any ordered field, we obtain
\[
(-1,0)=-(1,0)=-\mathbf{1}<\mathbf{0}.
\]

This contradicts the Trichotomy Property. 
\end{proof}